\documentclass[10pt,a4paper]{amsart}
\usepackage[margin=19mm]{geometry}
\usepackage{amsmath,amssymb,mathtools}
\usepackage[scr=rsfs,cal=euler]{mathalfa}
\usepackage{xfrac}
\usepackage[unicode,hidelinks,bookmarks=false]{hyperref}
\newcommand{\ie}{i.e.\ }
\newcommand{\cf}{cf.\ }
\newcommand{\resp}{respectively\ }
\newcommand{\Subsection}{\subsection}
\providecommand{\nobreakdash}{\nobreak\hbox{-}}
\setlength{\emergencystretch}{2em}
\multlinegap0pt
\usepackage{amssymb}
\usepackage{bm}
\usepackage{enumerate}
\usepackage{braket}
\usepackage{csquotes}
\usepackage{mathtools}
\usepackage{siunitx}
\usepackage{tikz}
\usepackage{float}
\newtheorem{lemma}{Lemma}
\newcommand{\A}{\mathcal{A}}
\newcommand{\R}{\mathbb{R}}
\newcommand{\s}{\mathbb{S}}
\title{A regularity result for $BV^{\A}(\Omega)$}
\author{Jakob Deutsch, Samuele Ricc\`o}
\date{Source: arXiv:2605.19504v1 (CC BY 4.0)}
\begin{document}
\maketitle

\noindent\emph{Source and licence.} A regularity result for $BV^{\A}(\Omega)$. Version arXiv:2605.19504v1 (CC BY 4.0), distributed by arXiv under the Creative Commons Attribution 4.0 International licence, http://creativecommons.org/licenses/by/4.0/, as recorded in the arXiv licence metadata for this exact version and shown on the abstract page https://arxiv.org/abs/2605.19504v1. Full licence notice: \texttt{licenses/arxiv-2605.19504-CC-BY-4.0.txt}.\par\smallskip

\noindent\textit{Editorial scope.} Counted unit: the article's lemma lem:translation\_lemma (source0.tex lines 462--474). The article's complete original proof of it is delivered with it (source0.tex lines 476--527, one \texttt{proof} environment of 52 lines). No mathematics is written, repaired or re-derived: the statement and the proof are the article's own lines, in the article's own notation, and no retained line is substituted or re-flowed.  Retained uncounted as the source-local context the proof needs: lines 310--313; lines 416--422; lines 448--450.  Nothing was omitted from any retained span, so the package carries no omission record.  No retained line contains a citation, so the wrapper carries no bibliography and no bibliography entry.  No retained line contains a figure environment, an image-inclusion command or drawing code, so no image is delivered and the package declares no asset dependency.  Editorial formatting only, in addition: an AMS \texttt{amsart} preamble that restates the article's own declarations and macros used by the retained lines, namely the environments \texttt{lemma} on separate counters, together with the standard packages those lines need.  The retained environments print in this document as Lemma 1 in order of appearance, whereas the article prints the same items with the numbering of its own full document.  The complete native source of the article as distributed in the arXiv e-print for this exact version is bundled unchanged as \texttt{sources/200919/source0.tex}.
\begin{equation}
\label{def:oper_first_ord}
    \A := \sum_{i = 1}^n A_i \partial_i
\end{equation}

\begin{lemma}\label{lem:polarization}
    Let $\A$, defined as in \eqref{def:oper_first_ord}, be elliptic and satisfy the rank-one property. Suppose that $(\xi, e), (\eta, f) \in \partial \sigma(\A)$. Then, there exists $\gamma \in \R$ with $\gamma \neq 0$ such that 
    \[
        (\xi + \lambda \eta, e + \gamma \lambda f) \in \partial \sigma (\A)
    \]
    for every $\lambda \in \R$.
\end{lemma}

\begin{lemma}\label{lem:directional_spectrum_surjectivity_projections}
    Let $\A$, defined as in \eqref{def:oper_first_ord}, be elliptic and satisfy the rank-one property. We have $P_{\R^n}(\partial \sigma(\A)) = \R^n$ and $P_{V^*}(\partial \sigma(\A)) = V^*$.
\end{lemma}

\begin{lemma}
\label{lem:translation_lemma}
    Let $\A$, defined as in \eqref{def:oper_first_ord}, be elliptic and satisfy the rank-one property. Let $K > 0$ and let $u \in L^1(\R^n; V)$ be such that 
    $$\| u \|_{L^1(\R^n; V)} < K.$$ 
    Suppose that for every $(\xi, v^*) \in \partial \sigma(\A) \cap (\s^{n-1} \times V^*)$ with $|v^*| = 1$ there exists $\omega_{(\xi, v^*)}: \R^+ \to \R^+$ non-decreasing with $\lim_{h \to 0} \omega_{(\xi, v^*)}(h) = 0$ such that for every $h \in (0,1)$ we have
    \[
        \int_{\R^n} |\langle v^*, u(x + h \xi)\rangle - \langle v^*, u(x) \rangle| \, dx \leq \omega_{(\xi, v^*)}(h).
    \]
    Then, there exist $h_0 \in (0,1)$ and $\tilde \omega : \R^+ \to \R^+$ non-decreasing with $\lim_{h \to 0} \tilde \omega(h) = 0$ (only depending on $K$ and finitely many $\omega_{(\xi, v^*)}$) such that for every $h \in (0, h_0)$ we have
    \begin{align}\label{eq:to_show}
        \int_{\R^n} |u(x + h \xi)- u(x)| \, dx \leq \tilde \omega(h).
    \end{align}
\end{lemma}

\begin{proof}
    Let $\xi \in \s^{n-1}$. Assume first that  
    \[
        \int_{\R^n} |u(x \pm h e_i)- u(x)| \, dx \leq \tilde \omega_i(|h|)
    \]
    holds for every $i = 1,\ldots, n$ and every $h \in (0,1)$ for some modulus of continuity $\tilde \omega_i$.
    Then,
    \begin{align}
    \label{eq:to_show_2}
        \int_{\R^n} |u(x + h \xi) - u(x)| \, dx 
        &\leq \sum_{j = 1}^N \int_{\R^n} \left|u \left(x + h\sum_{i = 1}^j \xi_i e_i \right) - u \left(x + h\sum_{i = 1}^{j -1} \xi_i e_i \right) \right| \, dx \nonumber \\
        &= \sum_{j = 1}^N \int_{\R^n} |u(x + h\xi_j e_j) - u(x)| \, dx  \nonumber \\
        &\leq \sum_{j = 1}^N \tilde \omega_j(h|\xi_j|)  \nonumber \\
        &\leq \sum_{j = 1}^N \tilde \omega_j(h).
    \end{align}
    It therefore suffices to only check \eqref{eq:to_show} for a finite number of $\xi$.
    \\
    Let $\xi \in \s^{n-1}$. Choose $v^* \in V^*$ such that $|v^*| = 1$ and $(\xi, v^*) \in \partial \sigma (\A)$, which is possible due to Lemma \ref{lem:directional_spectrum_surjectivity_projections}. Moreover, choose $(w_i^*)_{i = 1}^{N-1}$ such that $v^*, w_1^*, ..., w_{N - 1}^*$ form an orthonormal basis of $V^*$. Now, we again apply Lemma \ref{lem:directional_spectrum_surjectivity_projections} to get $\xi_{i} \in \s^{n-1}$ such that $(\xi_i, w_i^*) \in \partial \sigma(\A)$. Then,
    \[
        |u(x + h\xi) - u(x)| \leq |\langle v^*, u(x + h\xi) - u(x) \rangle| + \sum_{i = 1}^{N - 1}|\langle w_i^*,u(x + h\xi) - u(x)\rangle |.
    \]
    We directly estimate 
    \begin{align*}
        &|\langle w_i^*, u(x + h\xi) - u(x)\rangle| \\
        \leq ~ & |\langle w_i^*, u(x + h\xi) - u(x + \sqrt h \xi_{i})\rangle| + |\langle w_i^*, u(x + \sqrt h \xi_{i}) - u(x)\rangle|.
    \end{align*}
    Now, define $\zeta_i \in \s^{n-1}$ and $s_h^i \in \R$ for every $i = 1, \ldots, N-1$ via the equation
    \[
        x + h\xi = x + \sqrt h \xi_{i} + s_h^i \zeta_{i}
    \]
    or, equivalently,
    \[
        \zeta_i = (\xi_i + \sqrt h \xi)\frac{\sqrt h}{s_h^i}.
    \]
    By Lemma \ref{lem:polarization} for every $i = 1, \ldots, N-1$ there exists $\gamma_i \in \R$ such that $g_i^* := (w_i^* + \gamma_i \sqrt h v^*)/|w_i^* + \gamma_i \sqrt h v^*|$ fulfills $(\zeta_i, g_i^*) \in \partial \sigma(\A)$. We now notice that $|s_h^i| \leq 2\sqrt h$ from which we further infer $|\zeta_i - \xi_i| \leq 2\sqrt h$. By construction, we also have $|w^*_i - g^*_i| \leq L|\gamma_i||\sqrt h|$ for a constant $L > 0$ and small $h > 0$. From this, we now derive  
    \begin{align*}
        &|\langle w_i^* - g_i^*, u(x + h\xi)\rangle| + |\langle g_i^*, u(x + h\xi) - u(x + \sqrt h \xi_{i})\rangle| + |\langle w_i^* - g_i^*, u(x + \sqrt h \xi_{i})\rangle| \\
        \leq \ & (|u(x + h\xi)| + |u(x + \sqrt h\xi_i)|)|\gamma_i| \sqrt h + |\langle g_i^*, u(x + \sqrt h \xi + s_h^i \zeta_i) - u(x + \sqrt h \xi) \rangle|.
    \end{align*}
    Integrating gives 
    \begin{align*}
        &\int_{\R^n} |u(x + h\xi) - u(x)| \, dx \\
        & \hspace{4em} \leq \sum_{i = 1}^N \int_{\R^n} |\langle g_i^*, u(x + \sqrt h \xi + s_h^i \zeta_i) - u(x + \sqrt h \xi) \rangle | \, dx + 2 K \left(\sum_{i = 1}^N|\gamma_i| \right)\sqrt h \\
        & \hspace{4em} = \sum_{i = 1}^N \int_{\R^n} |\langle g_i^*, u(x + s_h^i \zeta_i) - u(x) \rangle | \, dx + 2 K \left(\sum_{i = 1}^N|\gamma_i| \right)\sqrt h \\
        & \hspace{4em} \leq \sum_{i = 1}^N \omega_{(\zeta_i, g_i^*)}(h) + 2 K \left(\sum_{i = 1}^N|\gamma_i| \right)\sqrt h. \\
    \end{align*}
    This yields \eqref{eq:to_show} with $\omega_\xi(h) := \sum_{i = 1}^N \omega_{(\zeta_i, g_i^*)}(h) + 2 K \left(\sum_{i = 1}^N|\gamma_i| \right)\sqrt h$ for an arbitrary $\xi \in \s^{n-1}$. Defining 
    \[
        \tilde \omega := \sum_{i = 1}^N \tilde \omega_{e_i}(h),
    \]
    we infer \eqref{eq:to_show} for general $\xi \in \s^{n-1}$ from \eqref{eq:to_show_2}.
\end{proof}

\end{document}
